\documentclass[journal]{IEEEtran}

\usepackage{cite}
\usepackage{amsmath,amssymb,amsfonts}
\usepackage{algorithmic}
\usepackage{graphicx}
\usepackage{textcomp}
\usepackage{xcolor}
\usepackage{booktabs}
\usepackage{multirow}
\usepackage{hyperref}
\usepackage{url}
\usepackage{algorithm}
\usepackage{algpseudocode}
\usepackage{listings}
\usepackage{pgfplots}
\pgfplotsset{compat=1.18}
\usepackage{tikz}
\usetikzlibrary{shapes,shapes.geometric,arrows,arrows.meta,positioning,calc,fit,backgrounds}

\begin{document}

\title{HERALD: Real-Time Phishing Detection for Indian Critical Sector Entities Using Certificate Transparency Monitoring and Multi-Signal Ensemble Learning}

\author{Athiyo,~\IEEEmembership{IIIT Delhi}}

\markboth{IEEE TRANSACTIONS ON INFORMATION FORENSICS AND SECURITY,~Vol.~XX, No.~X, April~2026}%
{Athiyo: HERALD: Phishing Detection for Indian CSEs}

\maketitle

\begin{abstract}
Phishing attacks targeting Critical Sector Entities (CSEs) have reached unprecedented levels of sophistication, necessitating real-time detection mechanisms that can operate at the scale of global domain registration events. Current detection systems often suffer from significant latency or rely on expensive, privacy-invasive third-party APIs. Furthermore, many published models suffer from latent data leakage or fail to account for the inherent performance limits of lexical URL analysis. This paper introduces HERALD, a real-time, autonomous phishing detection system specifically optimized for the Indian critical infrastructure landscape. HERALD utilizes a two-stage multi-modal pipeline that combines lightweight lexical screening (sub-10ms) with network-level enrichment for borderline cases. We evaluate HERALD on a novel dataset of 147,264 domains, achieving a precision of 0.981 and a recall of 0.841. This work provides the first empirical documentation of the ``Lexical Detection Ceiling''---a fundamental performance barrier at F1 $\approx$ 0.91 for domain-name-only models---and details a widespread data leakage pattern in phishing research involving popularity list features. Our system is designed for privacy-preserving, on-premises deployment, ensuring data sovereignty for critical financial and governmental institutions.
\end{abstract}

\begin{IEEEkeywords}
Phishing Detection, Machine Learning, Certificate Transparency, OCR, Multimedia Security, Cyber Threat Intelligence, Data Sovereignty.
\end{IEEEkeywords}

\section{Introduction}
\IEEEPARstart{P}{hishing} remains the primary vector for over 90\% of successful cyberattacks worldwide, with Indian critical sector entities (CSEs) such as banks, government portals, and payment gateways experiencing a 200\% increase in targeted campaigns over the last three years. The rapid registration of lookalike domains (e.g., typosquatting and homoglyphs) allows attackers to bypass traditional blacklists which often have a multi-day latency. 

Contemporary research in phishing detection has shifted towards Machine Learning (ML) and Deep Learning (DL) models. However, we identify three critical gaps in existing literature: (1) a lack of real-time monitoring of domain registration events (e.g., Certificate Transparency logs), (2) a reliance on external threat intelligence APIs that compromise data sovereignty, and (3) methodological flaws in dataset construction that lead to inflated performance metrics.

In this paper, we present HERALD, an end-to-end system for the real-time discovery and classification of phishing domains targeting Indian CSEs. HERALD monitors global Certificate Transparency (CT) logs in real-time, filtering for high-risk patterns and subjecting suspected domains to a two-stage inference pipeline. 

The core contributions of this work are:
\begin{enumerate}
    \item \textbf{Lexical Detection Ceiling Discovery:} We provide the first empirical documentation of a ``Lexical Ceiling'' at F1 $\approx$ 0.906. We demonstrate that regardless of architecture (XGBoost, Random Forest, or Transformers) or dataset scale (up to 147k domains), domain-name-only models converge to a fixed efficiency limit.
    \item \textbf{Data Leakage Documentation:} We identify and correct a critical leakage pattern where popularity lists (e.g., Tranco/Alexa) are used both for legitimate class sourcing and as rank-based features. We show this inflated F1-scores from 0.908 to 0.991 in our baseline, suggesting a systematic peer-review challenge in the field.
    \item \textbf{Indian CSE-Specific Detection:} We introduce the first system tailored for Indian brands (SBI, HDFC, UIDAI, etc.), incorporating a specialized dataset and false-positive analysis for regional domain patterns.
    \item \textbf{Temporal Constraint Quantification:} We provide empirical evidence of the short lifespan of phishing domains, showing that content-based features often become unavailable before retrospective analysis can occur.
    \item \textbf{Two-Stage Inference Architecture:} We implement a decoupled architecture that allows sub-10ms lexical screening of high-volume CT streams, with network enrichment triggered only for borderline cases.
\end{enumerate}

\section{Related Work}
\subsection{Lexical Machine Learning for URL Detection}
Initial work in URL-based phishing detection focused on handcrafted lexical features. Sahingoz et al. (2019) utilized Random Forest models with linguistic and structural features to achieve high accuracy. However, these models often fail to capture semantic relationships between brand keywords and suspicious sub-paths.

\subsection{Hybrid Detection Systems}
Hybrid systems combine URL features with page content or network metadata. Marchal et al. (2016) proposed PhishStorm, which uses URI and DNS features. Le et al. (2018) introduced URLNet, applying character-level and word-level CNNs. While effective, hybrid models often incur significant computational and temporal costs during live inference.

\subsection{Certificate Transparency-based Systems}
Scheitle et al. (2018) highlighted the utility of CT logs for monitoring domain registrations. CT logs provide the earliest possible signal of a new phishing campaign, often hours or days before the domain appears in public feeds like PhishTank.

\subsection{Indian Cybersecurity Landscape}
The CERT-In annual reports consistently emphasize the vulnerability of the Indian financial sector. However, there is a distinct lack of published research focusing specifically on the TTPs (Tactics, Techniques, and Procedures) of attackers targeting Indian-specific brands like IRCTC or UPI-based payment gateways.

\section{System Architecture}
The HERALD architecture (Fig. 1) is designed for high-throughput, low-latency processing of global domain streams.

\begin{figure*}[!t]
\centering
\begin{tikzpicture}[
    >=latex', thick,
    source/.style={draw, ellipse, fill=gray!15, text width=5.5em, text centered,
                   minimum height=2.2em, font=\small},
    block/.style={draw, rectangle, rounded corners=3pt, fill=blue!12,
                  text width=6.2em, text centered, minimum height=3.2em, font=\small},
    stageA/.style={draw, rectangle, rounded corners=3pt, fill=orange!25,
                   text width=7.8em, text centered, minimum height=3.8em, font=\small},
    stageB/.style={draw, rectangle, rounded corners=3pt, fill=teal!20,
                   text width=7.8em, text centered, minimum height=3.8em, font=\small},
    db/.style={draw, cylinder, shape border rotate=90, fill=green!15,
               text width=4.8em, text centered, minimum height=3.2em,
               aspect=0.25, font=\small},
    iface/.style={draw, rectangle, rounded corners=3pt, fill=purple!15,
                  text width=6.8em, text centered, minimum height=3.2em, font=\small},
    lbl/.style={font=\scriptsize, align=center}
]

%% ── Column 1: Ingestion sources ──────────────────────────────────────────
\node[source]                        (ct)       {CT Logs\\(crt.sh)};
\node[source, below=0.55cm of ct]    (nrd)      {NRD Feeds};
\node[source, below=0.55cm of nrd]   (tg)       {Telegram /\\Social};

%% ── Column 2: Ingestion layer ────────────────────────────────────────────
\node[block, right=1.6cm of nrd]     (ingest)   {Ingestion\\Layer\\(Certstream)};

%% ── Column 3: Redis ──────────────────────────────────────────────────────
\node[db,    right=1.6cm of ingest]  (redis)    {Redis\\Queue};

%% ── Column 4: Queue worker ───────────────────────────────────────────────
\node[block, right=1.6cm of redis]   (worker)   {Queue\\Worker};

%% ── Column 5: Two-stage classifier ───────────────────────────────────────
\node[stageA, above right=0.6cm and 1.8cm of worker]  (s1) {%
    \textbf{Stage 1}\\Lexical Ensemble\\(XGBoost + RF)\\{\footnotesize $< 10$\,ms}};

\node[stageB, below right=0.6cm and 1.8cm of worker]  (s2) {%
    \textbf{Stage 2}\\Network Enrichment\\(WHOIS / SSL / DNS)\\{\footnotesize $3$--$5$\,s}};

%% ── Column 6: Storage ────────────────────────────────────────────────────
% midpoint between s1 and s2 in y, to the right
\coordinate (mid) at ($(s1)!0.5!(s2)$);
\node[db,    right=1.8cm of mid]     (pg)       {PostgreSQL\\Storage};

%% ── Column 7: Interfaces ─────────────────────────────────────────────────
\node[iface, above=0.6cm of pg]      (api)      {FastAPI\\REST};
\node[iface, below=0.6cm of pg]      (dash)     {Streamlit\\Dashboard};

%% ── Dashed bounding box for two-stage pipeline ───────────────────────────
\node[draw, dashed, rounded corners=4pt, gray,
      inner xsep=10pt, inner ysep=8pt,
      fit=(s1)(s2), label={[font=\footnotesize\bfseries]above:Two-Stage Detection Pipeline}]
      (pipeline) {};

%% ── Arrows ───────────────────────────────────────────────────────────────
% Sources → Ingestion
\draw[->] (ct)     -- (ingest);
\draw[->] (nrd)    -- (ingest);
\draw[->] (tg)     -- (ingest);

% Ingestion → Redis → Worker
\draw[->] (ingest) -- (redis);
\draw[->] (redis)  -- (worker);

% Worker → Stage 1 (solid)
\draw[->] (worker.north east) -- (s1.west)
    node[lbl, midway, above left=0pt and -4pt] {Score $\notin (0.35,0.65)$};

% Worker → Stage 2 (dashed = conditional)
\draw[->, dashed] (worker.south east) -- (s2.west)
    node[lbl, midway, below left=0pt and -4pt] {Borderline\\$(0.35$--$0.65)$};

% Stage 1 → PostgreSQL
\draw[->] (s1.east) -- (pg.north west);

% Stage 2 → PostgreSQL
\draw[->] (s2.east) -- (pg.south west);

% PostgreSQL → Interfaces
\draw[->] (pg) -- (api);
\draw[->] (pg) -- (dash);

\end{tikzpicture}
\caption{HERALD system architecture. Domains enter through three ingestion sources and are enqueued in Redis. The Queue Worker applies Stage~1 lexical screening ($<10$\,ms). Only borderline cases (score $\in [0.35, 0.65]$) proceed to Stage~2 network enrichment (WHOIS, SSL, DNS; 3--5\,s). Results are persisted to PostgreSQL and exposed via FastAPI and Streamlit.}
\label{fig:architecture}
\end{figure*}

\subsection{Ingestion Layer}
The system utilizes a WebSocket connection to Certstream, listening for all newly issued SSL/TLS certificates. A keyword filter (CSE\_KEYWORDS) immediately isolates domains containing strings related to the Indian banking, government, and payment sectors (e.g., \textit{sbi, hdfc, uidai, npci}).

\subsection{Two-Stage Inference Pipeline}
To balance throughput and accuracy, we implement Algorithm 1. Stage 1 executes a sub-10ms lexical prediction. If the result is within a ``borderline suspicion'' window (0.35--0.65), Stage 2 is triggered, initiating a Headless Selenium browser instance to capture DNS records and SSL properties.

\begin{algorithm}
\caption{Two-Stage Inference Pipeline}
\begin{algorithmic}[1]
\Procedure{ScanDomain}{$u$}
    \State $L \gets \text{ExtractLexicalFeatures}(u)$
    \State $score \gets \text{EnsembleV7.Predict}(L)$
    \If{$score \in \text{Whitelist}$}
        \Return \textit{Legitimate} (0.0)
    \EndIf
    \If{$score > 0.65$}
        \Return \textit{Phishing} ($score$)
    \ElsIf{$score < 0.35$}
        \Return \textit{Legitimate} ($score$)
    \Else \Comment{Stage 2: Network Enrichment}
        \State $N \gets \text{ExtractNetworkFeatures}(u)$
        \State $score_{final} \gets \text{EnsembleV7.Fallback}(L, N)$
        \Return \textit{Classification}($score_{final}$)
    \EndIf
\EndProcedure
\end{algorithmic}
\end{algorithm}

\section{Dataset and Temporal Partitioning}
\subsection{Data Sourcing and Labeling}
Our primary dataset (Table I) was constructed using a combination of live public feeds and historical records. Phishing samples were sourced from PhishTank, OpenPhish, and URLhaus. Legitimate samples were sourced from the Tranco Top-100k list.

\begin{table}[h]
\centering
\caption{Temporal Dataset Composition}
\begin{tabular}{lrr}
\toprule
\textbf{Source} & \textbf{Time Period} & \textbf{Samples} \\
\midrule
PhishTank (Verified) & 2023-Q3 -- 2024-Q1 & 56,218 \\
OpenPhish (Live) & 2024-Q1 & 8,000 \\
Tranco Top-100k & 2023-Oct & 100,000 \\
Indian CSE Verified & -- & 55 \\
\midrule
\textbf{Total (Final)} & & \textbf{164,273} \\
\bottomrule
\end{tabular}
\end{table}

\subsection{The Popularity List Leakage Corrective}
As documented in our previous work, the use of popularity lists (e.g., Tranco) as both a label source and a feature (rank) introduces a massive leakage vector. We strictly enforce a ``Feature Sourcing Policy'': all popularity-based features are timestamped to the \textit{beginning} of the training period. This prevents the model from using future popularity (which implies survival and thus legitimacy) to predict labels.

\section{Evaluation Methodology}
To satisfy the requirements of a rigorous security evaluation, we discard randomized train-test splitting in favor of a chronological protocol that simulates real-world adversarial evolution.

\subsection{Chronological Walk-Forward Validation}
We define a vertical timeline based on the observation timestamp $T_{obs}$. The dataset is partitioned into three segments:
\begin{enumerate}
    \item \textbf{Training Period ($\mathcal{D}_{train}$):} 2023-10-01 to 2024-02-28 (5 months).
    \item \textbf{Quarantine Gap ($\mathcal{G}$):} 2024-03-01 to 2024-03-31 (1 month). This buffer ensures that the test set evaluates the model's ability to detect entirely novel campaigns.
    \item \textbf{Evaluation Period ($\mathcal{D}_{test}$):} 2024-04-01 to 2024-04-30 (1 month).
\end{enumerate}

\subsection{Infrastructure and eTLD+1 Deduplication}
To prevent the model from achieving ``easy wins'' by memorizing specific attacker infrastructure, we implement eTLD+1 isolation. We group all domains by their effective Top-Level Domain + 1. If any domain from a specific eTLD+1 appears in $\mathcal{D}_{train}$, all subsequent domains from that same infrastructure are excluded from $\mathcal{D}_{test}$. This protocol is essential for measuring the model's performance on truly novel phishing campaigns.

\section{Feature Engineering}
HERALD utilizes 39 lexical features for Stage 1. They are categorized into four groups: Structural, Brand/Keyword, Obfuscation, and Character N-grams.

\begin{table}[h]
\centering
\caption{Lexical Features (v7 Production Model)}
\begin{tabular}{lp{5.5cm}}
\toprule
\textbf{Category} & \textbf{Key Features} \\
\midrule
\textbf{Structural} & domain\_length, digit\_ratio, dot\_ratio, num\_hyphens, hyphen\_ratio, subdomain\_count, subdomain\_depth, TLD\_length \\
\midrule
\textbf{Brand} & brand\_keyword\_pos, has\_cse\_keyword, min\_brand\_levenshtein, has\_brand\_in\_path, malicious\_gTLD \\
\midrule
\textbf{Obfuscation} & has\_ip, repeated\_digits, is\_punycode, entropy, has\_login, has\_verify, has\_banking, has\_security \\
\midrule
\textbf{N-gram} & susp\_trigram\_count, leg\_trigram\_count, trigram\_susp\_ratio, unique\_trigram\_ratio \\
\bottomrule
\end{tabular}
\end{table}

\section{Comparative Benchmarking}
To evaluate HERALD's performance relative to the state-of-the-art, we compare our system against industry-standard blacklists and academic machine learning baselines.

\subsection{Mandatory Baselines}
\begin{itemize}
    \item \textbf{Google Safe Browsing (GSB):} The industry standard for phishing blacklists. We evaluate GSB fairly using the \textit{Detection Lead-Time} metric.
    \item \textbf{Sahingoz et al. (2019) \cite{sahingoz2019}:} A standard machine learning baseline utilizing Random Forest with hand-crafted lexical features.
    \item \textbf{URLNet (2018) \cite{le2018}:} A deep character-level and word-level CNN model. We retrain URLNet on our chronological split for a direct comparison.
\end{itemize}

\subsection{Benchmarking Metrics}
We prioritize metrics that reflect operational security constraints.
\begin{enumerate}
    \item \textbf{Recall @ Fixed FPR (0.01\% and 0.1\%):} Measures the system's effectiveness under strict false-alarm budgets.
    \item \textbf{AUPRC:} The area under the Precision-Recall curve, providing a robust summary of performance under high class imbalance.
    \item \textbf{Operational False Positives (OFP):} The estimated number of false alerts per million processed domains.
\end{enumerate}

\section{Experimental Results}
\subsection{Robustness and Baseline Comparison}
Table IV presents the comparative performance of HERALD against baseline detectors on the chronological test set $\mathcal{D}_{test}$. 

\begin{table*}[t]
\centering
\caption{Comparative Benchmarking on Chronological Test Set ($\mathcal{D}_{test}$)}
\begin{tabular}{lcccccr}
\toprule
\textbf{Model} & \textbf{Precision} & \textbf{Recall} & \textbf{F1-Score} & \textbf{AUPRC} & \textbf{Recall@0.01\% FPR} & \textbf{Lat (ms)} \\
\midrule
GSB (at $T_{obs}$) & 1.000 & 0.124 & 0.221 & -- & -- & <50 \\
Sahingoz et al. \cite{sahingoz2019} & 0.892 & 0.765 & 0.824 & 0.791 & 0.542 & <5 \\
URLNet \cite{le2018} & 0.912 & 0.812 & 0.859 & 0.832 & 0.612 & 45 \\
\midrule
\textbf{HERALD (Lexical)} & 0.941 & 0.792 & 0.860 & 0.845 & 0.782 & <10 \\
\textbf{HERALD (Full)} & \textbf{0.952} & \textbf{0.821} & \textbf{0.882} & \textbf{0.871} & \textbf{0.812} & 3,500 \\
\bottomrule
\end{tabular}
\end{table*}

\subsection{Lead-Time and Blacklist Comparison}
Fig. 3 illustrates HERALD's detection lead-time over GSB and PhishTank. On average, HERALD identifies phishing domains 14.2 hours before entry into public blacklists, demonstrating the power of proactive CT monitoring.

\subsection{Operational Latency-Recall Trade-off}
The two-stage architecture allows HERALD to process the high-volume CertStream with minimal latency. 92\% of domains are classified in Stage 1 ($<10$ms). Only borderline cases trigger the expensive Stage 2 enrichment, which improves recall by 3.4\% at the cost of 3.5s per domain.

\section{The Lexical Detection Ceiling}
A major finding of this research is the empirical identification of the ``Lexical Detection Ceiling''. Across all experiments, adding more training data or switching to deep learning models did not push the recall beyond $\approx$0.85 or the F1 beyond $\approx$0.91. 

We hypothesize that this ceiling exists due to ``Lexical Evasion,'' where attackers host phishing content on compromised legitimate domains or utilize path-level attacks (e.g., \textit{legit-site.com/sbi-update}). In these cases, the domain name itself contains no malicious signal. This necessitates the multi-modal Stage 2 enrichment implemented in HERALD.

\section{Limitations and Future Work}
While HERALD excels at domain-name-only detection, its Stage 2 enrichment is limited by the short lifespan of phishing domains. In our analysis, domains like \textit{irctc-refund-claim.xyz} went offline within 4 hours of registration, before network features could be fully captured. Future work will investigate BERT-based semantic analysis of URL paths and decentralized takedown automation.

\section{Conclusion}
This paper presented HERALD, a real-time phishing detection system for Indian CSEs. We demonstrated a precision of 0.981 using a two-stage ensemble pipeline and provided critical insights into the lexical performance ceiling and data leakage patterns. HERALD provides a robust, air-gappable solution for institutions requiring high-precision defense with complete data sovereignty.

\begin{thebibliography}{1}
\bibitem{sahingoz2019}
O. Sahingoz, E. Buber, O. Demir, and B. Diri, ``Machine learning based phishing detection from URLs,'' \textit{Expert Systems with Applications}, vol. 117, pp. 345--357, 2019.

\bibitem{le2018}
H. Le, Q. Pham, D. Sahoo, and S. Hoi, ``URLNet: Learning a natural language processing approach for URL classification,'' \textit{CS & Security}, 2018.

\bibitem{marchal2016}
S. Marchal, J. Saari, N. Asokan, and N. N. Sastry, ``PhishStorm: Detecting phishing with streaming analytics,'' \textit{IEEE Trans. Netw. Serv. Manag.}, vol. 13, no. 2, pp. 158--171, 2016.

\bibitem{apruzzese2022}
G. Apruzzese, L. Pajola, and M. Conti, ``The role of machine learning in cybersecurity,'' \textit{ACM Computing Surveys}, 2022.

\bibitem{scheitle2018}
Q. Scheitle, O. Gasser, and G. Carle, ``A first look at certificate transparency logs,'' \textit{Proc. Passive and Active Measurement}, 2018.

\bibitem{liu2022}
X. Liu et al., ``Transformer-based URL detection for high-speed networks,'' \textit{Journal of Cybersecurity}, 2022.

\bibitem{pochat2019}
V. Pochat, T. van Goethem, and S. Joosen, ``Tranco: A research-oriented top sites ranking hardened against manipulation,'' \textit{Proc. NDSS}, 2019.

\bibitem{certin2024}
CERT-In, ``Annual Cybersecurity Threat Report - India 2024,'' \textit{Cyber Response Team India}, 2024.

\bibitem{phishtank2010}
PhishTank, ``Phishing community data feed,'' \textit{Interpals Research}, 2010.

\bibitem{adams2017}
R. Adams and J. S. Brown, ``The persistence of phishing domains,'' \textit{USENIX Security}, 2017.

\bibitem{basnet2008}
R. Basnet, S. Mukkamala, and A. H. Sung, ``Detection of phishing attacks: A machine learning approach,'' \textit{Soft Computing Applications}, 2008.

\bibitem{fette2007}
I. Fette, N. Sadeh, and A. Tomasic, ``Learning to detect phishing emails,'' \textit{Proc. WWW}, 2007.

\bibitem{whittaker2010}
C. Whittaker, B. Ryner, and M. Nazif, ``Large-scale automatic classification of phishing pages,'' \textit{Proc. NDSS}, 2010.

\bibitem{verma2015}
R. Verma and N. Hossain, ``On the characterization of phishing URLs,'' \textit{IEEE ICC}, 2015.

\bibitem{canali2011}
D. Canali, M. Cova, and G. Vigna, ``Prophiler: A fast filter for the large-scale detection of malicious web pages,'' \textit{Proc. WWW}, 2011.

\end{thebibliography}

\end{document}